\input{../tex-inputs/latex-leseansicht-vorspann}

\section[Arthur und Olga Schnitzler an Gustav Schwarzkopf, {[}11. 1. 1904?{]}]{L04425 Arthur und Olga Schnitzler an Gustav Schwarzkopf, {[}11. 1. 1904?{]}}
\nopagebreak\mylabel{L04425v}
\rehead{ }\normalsize\beginnumbering\briefempfaengerindex{Schwarzkopf, Gustav@\textsc{Schwarzkopf, Gustav}!zzzSchnitzler, Olga@\emph{von Olga Schnitzler}!1904-01-111@{{[}11. 1. 1904?{]}}|(be}\briefempfaengerindex{Schwarzkopf, Gustav@\textsc{Schwarzkopf, Gustav}!zzzSchnitzler, Arthur@\emph{von Arthur Schnitzler}!1904-01-111@{{[}11. 1. 1904?{]}}|(be}
\toendnotes[C]{\smallbreak\pagebreak[2]}
\correspDesc{Versand  durch Arthur Schnitzler, Olga Schnitzler am [11. 1. 1904?] in Semmering
\newline{}Übermittlung  am [11. 1. 1904?] in Semmering
\newline{}Zustellung  am 13. [1. 1904?] in Semmering
\newline{}Erhalt  durch Gustav Schwarzkopf am 13. [1. 1904?] in Wien}\toendnotes[C]{\smallbreak}
\Standort{DLA, A:Schnitzler, HS.1985.1.1897.}
\physDesc{Bildpostkarte, 83 Zeichen
\newline{}Handschrift Arthur Schnitzler: Bleistift, deutsche Kurrent
\newline{}Handschrift Olga Schnitzler: Bleistift
\newline{}Versand: 1) Stempel: »\nobreak{}\oindex{Semmering@\textbf{Semmering}|pwk}Semmering\nobreak{}«.   2) Stempel: »\nobreak{}\oindex{I., Innere Stadt@\textbf{I., Innere Stadt}|pwk}Wien 1/1 1, 13 {[}1. 04{]}, 10½–12V., Bestellt\nobreak{}«. }\toendnotes[C]{\smallbreak}\pstart{}{\pb}Herrn Guſtav Schwarzkopf\pend{}\pstart{}Wien I\oindex{I., Innere Stadt@\textbf{I., Innere Stadt}|pw}\pend{}\pstart{}Tiefer Graben 23\oindex{Wien@\textbf{Wien}!I., Innere Stadt@\textbf{I., Innere Stadt}!Tiefer Graben 23@\textbf{Tiefer Graben 23}|pw}.\pend{}{\bigskip}
\pstart
           {\pb}\textcolor{gray}{\textbf{Semmering}}\oindex{Semmering@\textbf{Semmering}|pw}\hfill \textcolor{gray}{\textbf{Waldpromenade zur Meierei\oindex{Meierei des Südbahnhotels@\textbf{Meierei des Südbahnhotels}|pw}.}}\pend
           \vspace{1em}\pstart {\pb}\label{K_L04425-1v}\edtext{Herzliche Grüße}{\lemma{\textnormal{\emph{Herzliche Grüße}}}\Cendnote{\textnormal{Die Karte ist undatiert, enthält aber zwei Hinweise. Ein kleiner Druckvermerk weist auf »1904« als frühestes mögliches Jahr. 
                  Beim Poststempel des Empfangs ist die Tagesziffer »13« deutlich zu erkennen. Vergleicht man
                  alle Aufenthalte am Semmering\oindex{Semmering@\textbf{Semmering}|pwk} in dieser Zeit mit einer Anwesenheit
                  am 11. oder 12. eines Monats, so ergibt sich nur der 11. 1. 1904
                  oder 12. 1. 1904 als wahrscheinliches Datum. Das deckt sich damit, 
                  dass die Karte ein Wintermotiv zeigt und Schnitzler am 11. 1. 1904
                  in der Meierei\oindex{Meierei des Südbahnhotels@\textbf{Meierei des Südbahnhotels}|pwk} war, die im Motiv der Karte erwähnt wird.}}}\label{K_L04425-1}! Ihr \spacefill\mbox{ArthSchn}\pend{}\selectlanguage{ngerman}\vspace{1em}\pstart {[}hs. O. Schnitzler:{]} \spacefill\mbox{OlgaSch.}\pend{}\selectlanguage{ngerman}\endnumbering\briefempfaengerindex{Schwarzkopf, Gustav@\textsc{Schwarzkopf, Gustav}!zzzSchnitzler, Olga@\emph{von Olga Schnitzler}!1904-01-111@{{[}11. 1. 1904?{]}}|)be}\briefempfaengerindex{Schwarzkopf, Gustav@\textsc{Schwarzkopf, Gustav}!zzzSchnitzler, Arthur@\emph{von Arthur Schnitzler}!1904-01-111@{{[}11. 1. 1904?{]}}|)be}\mylabel{L04425h}
\begin{anhang}
\end{anhang}\newcommand{\dateiname}{L04425}\newcommand{\titel}{Arthur und Olga Schnitzler an Gustav Schwarzkopf, [11. 1. 1904?]}\newcommand{\editorInnen}{Herausgegeben von Jahnke, SelmaMüller, Martin Anton}\input{../tex-inputs/latex-leseansicht-abspann}
